


\poempart{Reaching Distant Peaks}{Breaking the Laws of the Physical World}{

This part is mainly about ``A World Apart'' and ``Shadow Clone Jutsu.'' Together, they answer how ``connection'' and ``replication'' form the infrastructure of the Fifth Dimension---that the essence of communication is not transmission but resonance, and that data needs depth rather than sheer volume.

\vspace{1em}

This part is light and fun. Through many fascinating stories from physics and information theory, it guides you to understand how the Fifth Dimension builds the link between the physical world and the digital world. Data---the most important means of production in the AI era---how it is generated, collected, and stored, all reexamined from the perspective of ``\textbf{completeness ratio}'' as a new way to think about big data.}

% \part{Breaking the Laws of the Physical World}
\chapter{A World Apart---Connectivity and Latency}

\begin{myquote}
Is the foundation of the world matter, energy, or information?\\
Matter occupies space, energy drives motion, yet information consumes only time.\\
In the Fifth Dimension, distance is compressed, time becomes the sole currency, and consensus is an expensive luxury.
\end{myquote}

\section{Electricity's Second Identity: When Energy Learned to Speak}

\subsection{The Wave-Particle Duality of Energy and Signal}

For a very long time, our understanding of electricity was confined to energy. We built generators to simulate the power of muscles---turning millstones, lifting heavy loads. Electricity in those days was raw and brute, an extension of Newtonian mechanics.

\vspace{1em}

After the telegraph appeared, electricity quietly acquired a second identity: \textbf{it was both energy and signal.} Morse code has a subtle quality: it does not transmit letters directly, but rather transmits rhythm. Pauses, short tones, long tones---these form three fundamental symbols:

\vspace{1em}

Silence, for the moment;

Sound, but brief;

Sound, but slightly longer.

\vspace{1em}

They are neither as rich and varied as language, nor as structureless as natural noise, but compressed into discrete events that can be counted. A skilled telegraph operator receives not a static pattern of dots and dashes, but a stream of temporal undulations.

In other words, time was encoded into information: long and short, intervals and pauses, composing a script that dances along the axis of time.

\textbf{M}orse code is like a bridge---one end connected to humanity's millennia-old traditions of language and writing, the other to our modern binary world. As long as we can reliably distinguish a handful of states, we can build an entire language and logic upon them.

In 1937, Claude Shannon's moment arrived. In that celebrated master's thesis\footnote{Claude E. Shannon, \textit{A Symbolic Analysis of Relay and Switching Circuits}, 1937}, he revealed the equation that would alter the fate of humanity:
\begin{equation*}
\text{Switch (On/Off)} \equiv \text{Truth (True/False)}
\end{equation*}

\textbf{``True'' and ``false''} in Boolean algebra correspond perfectly to \textbf{``on'' and ``off''} in an electric circuit. Electric current carries information at the same time as it carries energy---and this transmission occurs at nearly the speed of light.

\begin{importantbox}
We gained speed---or rather, we gained the illusion of instantaneous arrival.
\end{importantbox}

\subsection{The Speed-of-Light Journey That Goes Nowhere}

Here is a counterintuitive fact: electricity, the carrier that sustains modern civilization, exhibits a peculiar split personality. At the microscopic level, electrons drift through a conductor at a snail's pace---mere millimeters per second\footnote{The propagation speed of a signal through a conductor depends on the dielectric constant.}; yet at the macroscopic level, the electric field establishes itself at nearly the speed of light.

Do not imagine that electricity transmits the information we wrote. Instead, understand it this way: it enables the operator at the other end of the line to write down an identical copy. The Fifth Dimension is not true teleportation---it is a pattern being re-illuminated at a different location.

\vspace{1em}

Recall the Ship of Theseus paradox mentioned earlier:

If every plank of a ship is gradually replaced, is it still the same ship?

\vspace{1em}

Personally, I believe that as long as the two ships are built from the exact same blueprint---the same information---and that information has not been tampered with in any way, then the new ship is the original ship.

What truly needs to be preserved is the invariance of the information structure.


\textbf{The essence of ``A World Apart'' is rebuilding the same structure elsewhere, not transporting the same matter.}


\section{The Race Between Bullet and Light: When Effect Outruns Cause}

\subsection{Why We Crave Speed So Desperately}

Shannon later told us: the essence of information is the elimination of uncertainty.

And throughout human history, uncertainty has almost always meant fear---or profit.

\vspace{1em}

Before the age of the telegraph, the boundary of power ended at the horizon.

Claude Chappe built his optical semaphore towers during the French Revolution, but they were still hostage to weather and the limits of the human eye. It was an extraordinarily fragile communication network---a single fog bank could sever Paris from its borders.

\vspace{1em}

Then financial capital entered the contest for certainty.

From the London Exchange to the Paris Bourse, hundreds of kilometers of geographical distance were folded into an information time-gap. The legend of the Rothschild family is, at its core, a story about who could ``see'' the outcome of Waterloo first within that gap.

\vspace{1em}

For financial speculators, \textbf{wealth is not born from the sharing of information, but from its asymmetry.} They exploited telegraph wires that carried information at the speed of light---not to bring people closer, but to open a new dimension by front-running along the axis of time.

\vspace{1em}

This gave rise to the absurd yet very real state of affairs in today's high-frequency trading: to gain a nanosecond-level advantage, trading firms will bore through the Appalachian Mountains to lay fiber-optic cable. In this dimension, time is no longer a flowing river but is sliced into countless tiny fragments, each with a price tag.

\vspace{1em}

At the scale of high-frequency trading, \textbf{the future is merely the past that someone else hasn't received yet.}

\begin{thinkbox}
If speed itself is profit, what happens when it is pushed to its extreme?
\end{thinkbox}

\newpage
\subsection{The Ultimate Speed Limit of Causality: Why the Bullet Cannot Outrun Light}

Yet this pursuit of the instantaneous is strictly forbidden by physics.

To understand the rules of the Fifth Dimension, we must first revisit the fundamental laws of the physical world.

\vspace{1em}

Imagine a universe without relativity's constraints---a Newtonian world where velocities can stack without limit. But before that, consider the laws of our actual reality: you stand a thousand meters away and pull the trigger. The speed of light $c$ is absolutely constant and far faster than any bullet, so the sequence of events is strictly locked in place:

1. You pull the trigger (cause).

2. Light reaches my retina at maximum speed---I see smoke at the muzzle.

3. Then the bullet reaches my chest---I feel searing pain (effect).


\begin{importantbox}
First we see the cause, then we suffer the effect. This is the fundamental constraint of causal order.
\end{importantbox}

But suppose the speed of light were no longer the fixed upper limit, and instead obeyed simple velocity addition: when you pull the trigger, the light reflected off the bullet would add the bullet's own velocity, traveling toward me at $c$ (speed of light) $+$ $v$ (bullet speed). This means the light carrying the image of the bullet in flight would outpace the light carrying the image of the muzzle flash.

\vspace{1em}

I would witness a scene of pure horror: walking down the street, I first see a vivid close-up of the bullet bearing down on me (perhaps I even feel the searing pain first), and only then do I see you in the distance raising your arm, pulling the trigger, the muzzle erupting with \textbf{a flash of fire that should have been a warning}.

\vspace{1em}

\textbf{Cause and effect have been inverted.}

In my subjective world, the ending occurs before the beginning---at the very instant of my death, you have not yet decided whether to fire. If such a thing were allowed, the history of the universe would no longer unfold along a single arrow of time, but dissolve into a tangled ball of yarn.

\vspace{1em}

Therefore, the speed of light $c$ is not merely the upper limit of velocity---it is the guardian of causality.
The universe imposes a maximum speed limit precisely to prevent the effect from outrunning the cause.

\subsection{Latency: The Price We Pay for Order}

This thought experiment is a crucial key to understanding the Fifth Dimension---the digital world.

In distributed systems (the internet), a bit is both a signal and that ``bullet.'' If the transmission of information lacked rigid constraints---if a server in Shanghai received a ``payment successful'' packet (effect) before it received the ``user clicked pay'' packet (cause)---the entire financial system would collapse like that causality-inverted murder.

\vspace{1em}

In the physical world, the laws of physics enforce causal order through the speed-of-light limit.

In the digital world, engineers simulate this causality using clocks and consistency protocols.

To prevent a ``causal inversion'' catastrophe in the digital realm, we must accept a cost: \textbf{latency}.
In distributed systems, latency is often not a mere technical defect---it is the waiting time we must pay to maintain causal order and consistency.

\vspace{1em}

We must wait for light (information) to traverse that distance, wait for every observer to confirm the moment ``the trigger was pulled.''
Latency is the technical trade-off that causality demands in the Fifth Dimension.

\textbf{In systems that coordinate across nodes, attempting to eliminate latency makes stability very hard to maintain.}

\section{The Privilege of Waiting: When Deep Thought Becomes a Luxury}

Imagine: at this very millisecond we click ``Buy,'' while a server on the other side of the planet is processing ``Sold Out.'' Which comes first? There is no God's-eye view. The system can only promise: after a period of propagation and synchronization, all nodes will eventually reach consensus.

\vspace{1em}

For ``A World Apart,'' there is no true simultaneity.

Distributed systems are forced to compromise: they accept \textbf{eventual consistency.}

Truth is not an instantaneous state, but a process of convergence.

In such systems, we often live amid temporary inconsistency, waiting for the truth to arrive.

\begin{thinkbox}
In the AGI era, different social strata have different capacities to endure waiting. What changes will this bring?
\end{thinkbox}

\newpage
\section{The Chasm Between System 1 and System 2}

If the latency of the traditional internet is the latency of physical transmission,
then the latency of the AI era is the \textbf{latency of reasoning}.

\vspace{1em}

As Daniel Kahneman revealed in \textit{Thinking, Fast and Slow},
the human brain possesses two fundamentally different systems:

\textbf{System 1 (fast thinking)}: intuition, conditioned reflex, millisecond-level response.

\textbf{System 2 (slow thinking)}: logic, reasoning, planning, minute-level response.

\vspace{1em}

In the Fifth Dimension, this distinction is reshaping the architecture of power:

\textbf{Cheap ``real-time'' (System 1)}

This is the tap water of the future. Token generation is extremely fast, marginal costs keep falling, and it is well suited for handling the bulk of everyday communication. Its hallmark: speed---but when depth of thought is absent, quality and reliability tend to plateau at ``good enough'' rather than ``the best.''

It is like a person whose prefrontal cortex has been removed, capable only of probabilistic conditioned reflex. Ordinary people will live in millisecond-level responses, receiving instant gratification, cheap entertainment, and probabilistic truth---unable to see the full picture of cause and effect before eventual consistency is reached.

\vspace{1em}

\textbf{Expensive ``deep thought'' (System 2)}

This is the gold of the future, and precisely why \textbf{compute will be the hardest currency of tomorrow}.

A reasoning model deliberates longer before answering---self-refuting, searching for verification, constructing chains of causation. The reasoning time and compute cost required for a high-quality decision may vastly exceed those of the ``instant answer'' mode. Whoever can afford the longer wait and the higher compute cost gains easier access to this privilege of deep thought.

They pay exorbitant compute fees and endure greater latency, all to calculate every possible trajectory of the bullet before pulling the trigger. If we count the price of compute and the capacity to tolerate latency as part of societal wealth, then we can say plainly:

\begin{importantbox}
In the AGI era, deliberation may be a luxury only the wealthy can afford.
\end{importantbox}

% We used to think technology existed to eliminate latency.
% We have come to realize that in this world saturated with the illusion of speed,
% \textbf{taking the time to compute a correct future} is the most expensive luxury of all.

\newpage
\section{The Tower Keeper's Betrayal: Cracks in ``A World Apart''}

% \subsection{Redundancy: Doubling Down on Verbosity}

Every magic has its price. In thermodynamic terms, the total entropy of an isolated system tends to increase; maintaining information and order always exacts a cost.

\vspace{1em}

In real physical circuits, every transmission is a battle against the chaos of the universe.

Resistance, thermal noise, radiation---all conspire to drag the signal back toward disorder.

\begin{myquote}
    Shannon first acknowledged that the world itself is not clean, then set out to answer:\\
    In a noisy channel, at what maximum rate can we still reliably transmit information?
\end{myquote}


His answer was: \textbf{redundancy.}

To combat channel noise, we often need to add redundancy---

to make sure you hear me clearly, I often have to say the same thing three times.

From this emerged an entire theoretical edifice: channel capacity, redundant coding, error-correcting codes.

We can trade doubled verbosity for more reliable silence.

\vspace{1em}

In engineering practice, TCP's handshakes and checksum comparisons are the ``tolls'' we pay: extra bits and extra time, exchanged for a small measure of certainty against the march of entropy.

\subsection{Man-in-the-Middle Attacks and Seals}

Do you remember that scene from \textit{The Count of Monte Cristo}?

The Count (Edmond Dant\`{e}s) bribed the keeper of a semaphore tower, having him alter a few critical movements of the mechanical arms. A single subtle manipulation, and Spanish bonds plummeted---a vast fortune was transferred in an instant.
Alexandre Dumas could hardly have imagined that his tale of vengeance would precisely foretell a major crisis of the digital world more than a century later: the man-in-the-middle attack.

\vspace{1em}

``A World Apart'' has a fatal weakness: \textbf{it is responsible for ``transmission,'' not for ``truth.''}

When we convert our consciousness, wealth, and secrets into streams of bits and send them coursing through a global fiber-optic network, we are in effect entrusting our most precious possessions to countless invisible ``tower keepers''---routers, switches, listening stations along submarine cables.

\textbf{I}n the Fifth Dimension, if the problem of trust cannot be solved, then ``A World Apart'' is not a wormhole to the future, but a trapdoor to the abyss.
To reduce this uncertainty, we invented the HTTPS protocol---
an ``unforgeable seal'' inscribed at the entrance and exit of every passage through ``A World Apart.''

\vspace{1em}

Through mechanisms of encryption, certificates, and integrity verification, it strives to deliver three guarantees: minimizing eavesdropping, minimizing tampering, and verifying the identity of the other party. Only with these mathematical and engineering safeguards does the Fifth Dimension establish a relatively solid ground to stand on.

\subsection{Noise in the Mind: The Incorrigible Misunderstanding}
All of this, at its core, acknowledges a single truth: ``A World Apart'' is not perfect---over a sufficient distance, information will inevitably decay into \textbf{noise}.
This leads to an even more dangerous consequence:

\begin{myquote}The sentence you received is, with high probability, not the sentence I originally said.\end{myquote}

This is the cost that ``communication'' must bear.

Physical signals distort. Bits flip.

\vspace{1em}

To ensure that ``the sentence you received'' is truly ``the sentence I originally said,'' we must introduce error-correcting codes---a form of mathematical redundancy.

\vspace{1em}

\textbf{But the more terrifying noise is not in the cables---it is in our minds.}

Context, bias, emotion---these are sources of interference far harder to eliminate than thermal noise.

We invented error-correcting codes to fix flipped bits.

But we have yet to invent an algorithm that can fix the gap in understanding between two people.

\vspace{1em}

Perhaps complete and lossless understanding between human beings was always close to impossible.

\newpage
\section{A Knowing Smile: Fewer Words, Greater Resonance}

Having solved the problem of physical transmission (fiber optics), and the problem of logical causality (latency), we finally arrive at another ultimate form of communication: \textbf{compression and resonance.}
Shannon told us: in a noisy channel, we must fight entropy with redundancy---communication seems destined to be inefficient.

But in the Fifth Dimension, we found a wormhole that bypasses physical limits: the \textbf{shared model}.

\subsection{When Two Minds Share the Same Dictionary}

Imagine that both sender and receiver hold an identical codebook---one containing the logic of everything in the universe. The logic of transmission would be fundamentally transformed:

\textbf{Traditional mode}: I want to show you an eagle soaring against a sunset. I need to send an entire high-resolution image, perhaps even a video's worth of data.

\textbf{Model mode}: Because your model already contains the concepts of ``sunset'' and ``eagle,'' I need only send a few bytes of instruction (Prompt/Seed).

\vspace{1em}

Transmission has been compressed to its extreme.

The massive stream of information vanishes, replaced by minimalist ``indices'' and ``pointers.''

This compression yields not merely bandwidth savings, but a supremely efficient semantic resonance.

When both parties share the same high-dimensional model, we enter a dreamlike state of ``low bandwidth, high signal-to-noise ratio.''
Fewer words\footnote{Real language models use tokens rather than words as their basic unit. Curious readers are encouraged to ask a GPT model what a token is.} can carry far grander meaning.

\vspace{1em}

This explains why the highest form of social communication is the \textbf{meme and humor}.

A joke is funny because the teller supplies only the first half (the trigger), and the listener---drawing on the same cultural background and life experience stored in their brain (the shared model)---instantly completes the absurd logic of the second half.

\vspace{1em}
The pinnacle of communication is not making yourself clearer,

but rather, even when I say nothing at all,

you can decode, from the slightest twitch of my eyebrow, that my whole world is filled with love for you.

\vspace{1em}

\textbf{T}hat \textbf{knowing smile}---

that is the moment when two independent world-models complete a perfect ``handshake'' in the Fifth Dimension.

In the future, communication between AGIs will work the same way.

They will not send each other encyclopedias---they will only exchange ``seeds'' and ``difference vectors.''

Whoever possesses the more aligned model possesses a faster ``speed of thought.''

\subsection{The Faster the Transmission, the Slower the Understanding}

In conventional thinking, the speed of ``A World Apart'' is seen as a one-directional indicator of progress:

{\kaiti
Greater bandwidth;

Lower latency;

Higher concurrent connections.
}

\vspace{1em}

But in lived experience, we often feel a paradox:

{\kaiti
Information moves faster and faster, yet attention grows ever more fragmented;

Messages arrive in near real-time, yet the time needed to truly discuss one thing has not shortened;

Even with planet-scale broadcasting power, we often find ourselves trapped in echo chambers of misunderstanding and outrage.
}

\vspace{1em}

This is the dilemma of ``A World Apart'' in the Fifth Dimension:

When the cost of connecting information approaches zero, and the cost of generating information also approaches zero:
\begin{importantbox}
The cost of filtering noise and distinguishing truth from falsehood skyrockets.
\end{importantbox}

Shannon defined information as the degree to which uncertainty is eliminated---this works perfectly in engineering.

But at the semantic level, there is an easily overlooked point:

\textbf{How much uncertainty can be eliminated depends on how much ``prior knowledge'' we share.}

\vspace{1em}

If two people's world-models are highly aligned,
a single glance, a single meme, a single abbreviation can convey an extraordinarily complex state transition;
if their world-models diverge greatly,
even with infinite bandwidth, they can only rebuild the other's mental context
in the most painstakingly verbose fashion, layer by layer.

\vspace{1em}

\textbf{F}rom this perspective, Shannon's channel capacity (bandwidth, latency) is the hardware ceiling;

and I would like to propose a complementary view: the degree of overlap between world-models is the software ceiling.

It is the latter that is often the true bottleneck to genuine mutual understanding.

The hardware ceiling can be raised through technology; the software ceiling can only be raised through shared experience.

\subsection{World-Model Overlap and the Futility of Shallow Socializing}

We can roughly imagine the difficulty of mutual understanding as:

the number of bits required to compile, in the other person's brain, an interpreter capable of running this explanation.

\vspace{1em}

Thus, a truly deep relationship

is, in essence, an iterative alignment of world-models over time---not the exchange of more messages.

The root of many quarrels lies in this:

``The same interpreter hasn't been installed yet,'' but we call the other person as if they were a local program.

\vspace{1em}

And genuine socializing will shift from collaboration for survival back to its most ancient essence---

\textbf{gathering for the sake of resonance.}

\vspace{1em}

We no longer sit together merely to finish a report.

We sit together simply because we want to feel the warmth of another body, to see the unguarded truth that surfaces after a glass of wine, to confirm that in this world---one that may well be overtaken by silicon-based intelligence---there is still another imperfect carbon-based being who shares this experience with me.

\begin{thinkbox}
After AGI arrives, what kind of socializing will count as real socializing?
\end{thinkbox}

\newpage
\section{Coda: The Truth of ``A World Apart''}

We call it ``A World Apart,'' but now we know---it is no magic that conjures distance away. It has never truly brought anything closer.

\vspace{1em}

It has merely established an agreement between two places:

When a lamp is extinguished here, a lamp is lit there;

When someone weeps here, a heart breaks there.

\vspace{1em}

The Ship of Theseus gave us the answer:

What matters is not the identity of the material, but the continuity of the pattern.

\vspace{1em}

When we invoke the magic of ``A World Apart,''

``I'' am not this body, but this pattern;

``You'' are not that address, but the interpreter that understands me;

``We'' are not a physical gathering, but an overlap of models.

\vspace{1em}

This is the magic of resonance in the Fifth Dimension---it was never meant to make us faster.

It was meant to help us understand that \textbf{some things are worth the wait.}
